\documentclass[manuscript,anonymous,review]{acmart}
\settopmatter{printacmref=false,printfolios=true}
\setcopyright{none}
\renewcommand\footnotetextcopyrightpermission[1]{}
\acmConference[ICCRD 2027]{2027 19th International Conference on Computer Research and Development}{January 15--17, 2027}{Singapore}
\usepackage{booktabs,mathtools}
\setlength{\emergencystretch}{1em}
\renewcommand{\thetable}{S\arabic{table}}
\begin{document}
\title{HDA-NQ: Supplementary Component Results}
\begin{abstract}
This supplement reports the complete component experiments at both studied prediction lengths. It accompanies the focused 96-step long-term forecasting analysis in the main manuscript and retains the longer-horizon accuracy--detail tradeoffs.
\end{abstract}
\ccsdesc[500]{Computing methodologies~Machine learning}
\keywords{quantile forecasting, component analysis, reproducibility}
\maketitle
\section{Scope and Protocol}
The component study uses ETTh1, history length 336, prediction lengths 96 and 720, and seeds 2024, 2025, and 2026. The six modified models are retrained independently at both lengths; full-model reference checkpoints are shared with the primary comparison. All use the same data split, training-only scaler, optimizer policy, and validation checkpoint selection as the final model. Each comparison changes only its named component.

The main text presents the 96-step results as a focused illustration, chosen after inspecting the completed results. This presentation choice does not alter any training configuration or test metric. Tables~\ref{tab:full-ablation} and \ref{tab:detail-ablation} retain both horizons, including component changes that improve the longest-horizon loss. No uniform benefit across horizons is implied.

\begin{table}[t]
\centering\small
\caption{Complete ETTh1 component study; pinball mean $\pm$ sample SD and mean crossing percentage over three seeds.}
\label{tab:full-ablation}
\begin{tabular}{@{}lrrrr@{}}
\toprule
Variant & Pinball 96 & Crossing 96 (\%) & Pinball 720 & Crossing 720 (\%) \\
\midrule
Full HDA-NQ & $0.12966\pm0.00018$ & $0$ & $0.18683\pm0.00068$ & $0$ \\
Fixed offsets & $0.13461\pm0.00147$ & $0$ & $0.18842\pm0.00098$ & $0$ \\
Shared embedding & $0.13978\pm0.00065$ & $0$ & $0.18880\pm0.00201$ & $0$ \\
Dense attention & $0.14158\pm0.00104$ & $0$ & $0.18971\pm0.00295$ & $0$ \\
No correction & $0.13299\pm0.00073$ & $0$ & $0.18586\pm0.00090$ & $0$ \\
Fixed + no correction & $0.13452\pm0.00139$ & $0$ & $0.18783\pm0.00169$ & $0$ \\
Independent head & $0.12997\pm0.00114$ & $9.81\times10^{-4}$ & $0.18591\pm0.00159$ & $1.91\times10^{-5}$ \\
\bottomrule
\end{tabular}
\end{table}

\begin{table}[t]
\centering\small
\caption{Complete local-variation diagnostics, mean $\pm$ sample SD over three seeds. Event thresholds use training-prefix changes.}
\label{tab:detail-ablation}
\begin{tabular}{@{}lrrrr@{}}
\toprule
Variant & $\Delta$MAE 96 & $\Delta$MAE 720 & Event MAE 96 & Event MAE 720 \\
\midrule
Full HDA-NQ & $0.2466\pm0.0036$ & $0.2663\pm0.0012$ & $0.7147\pm0.0084$ & $0.8598\pm0.0097$ \\
Fixed offsets & $0.2413\pm0.0005$ & $0.2659\pm0.0012$ & $0.7160\pm0.0037$ & $0.8634\pm0.0078$ \\
Shared embedding & $0.2460\pm0.0011$ & $0.2661\pm0.0058$ & $0.7393\pm0.0118$ & $0.8615\pm0.0207$ \\
Dense attention & $0.2506\pm0.0067$ & $0.2712\pm0.0020$ & $0.7421\pm0.0025$ & $0.8667\pm0.0032$ \\
No correction & $0.2727\pm0.0026$ & $0.3087\pm0.0031$ & $0.7490\pm0.0041$ & $0.8405\pm0.0089$ \\
Fixed + no correction & $0.2665\pm0.0019$ & $0.2862\pm0.0014$ & $0.7488\pm0.0124$ & $0.8632\pm0.0241$ \\
Independent head & $0.2473\pm0.0016$ & $0.2644\pm0.0027$ & $0.7156\pm0.0111$ & $0.8532\pm0.0082$ \\
\bottomrule
\end{tabular}
\end{table}

\section{Interpretation}
Learned positions reduce mean pinball relative to fixed references by 3.67\% at $H=96$ and 0.85\% at $H=720$. The corresponding reductions relative to a shared horizon embedding are 7.24\% and 1.05\%. Removing the correction raises first-difference error by 10.57\% and 15.94\%, respectively, but slightly improves 720-step pinball. Thus, preserving local variation does not guarantee a lower quantile loss. Independent outputs exhibit rare nonzero crossings and a slightly lower 720-step loss; the ordered head supplies a structural guarantee rather than a uniform accuracy advantage.

The fixed-reference and correction removals also form a two-by-two comparison. Contributions can interact with the correction branch and prediction length, so separate configuration changes should not be interpreted as additive causal effects. Independent-head results are reported with pinball and crossing rate; interval scores are omitted because unordered endpoints need not define valid central intervals.

The accompanying results files contain the complete unrounded component statistics, including sample standard deviations.
\section{Complete Interval Statistics}
Table~\ref{tab:full-intervals} expands the compact coverage and interval-score comparison in the main manuscript. Only dataset--horizon settings with all five models and three seeds complete are included, matching the main comparison. The horizon sets are listed in the main interval table; these are averaged within each training seed before computing the three-seed mean and sample standard deviation.
\begin{table}[t]
\centering\small
\caption{Interval statistics for the same completed settings as the main table. Mean $\pm$ sample SD over three seeds after averaging horizons within seed. Bold follows the main interval table.}
\label{tab:full-intervals}
\begin{tabular*}{\textwidth}{@{\extracolsep{\fill}}llrrr@{}}
\toprule
Dataset & Model & PICP90 (\%) & PINAW90 & Norm. IS90 \\
\midrule
ETTh1 & HDA-NQ (Ours) & $85.09\pm0.30$ & $0.2971\pm0.0023$ & $\boldsymbol{0.4409\pm0.0002}$ \\
ETTh1 & DLinear-NQ & $88.96\pm0.03$ & $0.3147\pm0.0005$ & $0.4621\pm0.0004$ \\
ETTh1 & PatchTST-NQ & $79.06\pm2.48$ & $0.2789\pm0.0260$ & $0.5514\pm0.0164$ \\
ETTh1 & TimeMixer-NQ & $\boldsymbol{90.13\pm1.88}$ & $0.4199\pm0.0114$ & $0.5319\pm0.0151$ \\
ETTh1 & iTransformer-NQ & $82.74\pm1.16$ & $0.2759\pm0.0069$ & $0.4784\pm0.0099$ \\
ETTh2 & HDA-NQ (Ours) & $86.94\pm0.22$ & $0.1945\pm0.0013$ & $\boldsymbol{0.2933\pm0.0008}$ \\
ETTh2 & DLinear-NQ & $\boldsymbol{92.87\pm0.11}$ & $0.3148\pm0.0038$ & $0.3790\pm0.0038$ \\
ETTh2 & PatchTST-NQ & $78.09\pm0.61$ & $0.2015\pm0.0021$ & $0.4441\pm0.0101$ \\
ETTh2 & TimeMixer-NQ & $68.21\pm0.06$ & $0.1908\pm0.0087$ & $0.5557\pm0.0106$ \\
ETTh2 & iTransformer-NQ & $85.59\pm0.88$ & $0.2482\pm0.0090$ & $0.4030\pm0.0015$ \\
ETTm1 & HDA-NQ (Ours) & $86.82\pm0.01$ & $0.2336\pm0.0002$ & $\boldsymbol{0.3451\pm0.0011}$ \\
ETTm1 & DLinear-NQ & $\boldsymbol{88.87\pm0.01}$ & $0.2611\pm0.0003$ & $0.4268\pm0.0003$ \\
ETTm1 & PatchTST-NQ & $84.91\pm0.85$ & $0.2489\pm0.0128$ & $0.4346\pm0.0206$ \\
ETTm1 & TimeMixer-NQ & $75.22\pm0.87$ & $0.2042\pm0.0061$ & $0.5265\pm0.0062$ \\
ETTm1 & iTransformer-NQ & $86.15\pm0.08$ & $0.2457\pm0.0023$ & $0.4167\pm0.0007$ \\
ETTm2 & HDA-NQ (Ours) & $87.99\pm0.09$ & $0.1573\pm0.0009$ & $\boldsymbol{0.2363\pm0.0007}$ \\
ETTm2 & DLinear-NQ & $91.79\pm0.50$ & $0.2352\pm0.0050$ & $0.3035\pm0.0015$ \\
ETTm2 & PatchTST-NQ & $90.50\pm0.22$ & $0.2142\pm0.0020$ & $0.2935\pm0.0029$ \\
ETTm2 & TimeMixer-NQ & $82.12\pm0.83$ & $0.1699\pm0.0068$ & $0.3281\pm0.0127$ \\
ETTm2 & iTransformer-NQ & $\boldsymbol{90.18\pm0.37}$ & $0.2236\pm0.0033$ & $0.3144\pm0.0012$ \\
Electricity & HDA-NQ (Ours) & $88.91\pm0.42$ & $0.1947\pm0.0036$ & $\boldsymbol{0.2965\pm0.0008}$ \\
Electricity & DLinear-NQ & $\boldsymbol{90.38\pm0.16}$ & $0.2181\pm0.0006$ & $0.3215\pm0.0001$ \\
Electricity & PatchTST-NQ & $89.44\pm0.16$ & $0.2019\pm0.0004$ & $0.3126\pm0.0010$ \\
Electricity & TimeMixer-NQ & $85.43\pm0.52$ & $0.1837\pm0.0022$ & $0.3310\pm0.0042$ \\
Electricity & iTransformer-NQ & $91.02\pm0.15$ & $0.2126\pm0.0030$ & $0.3111\pm0.0026$ \\
\bottomrule
\end{tabular*}
\end{table}

\end{document}
